\chapter{行列式}
    % \section{右下要素が$0$の対称行列の行列式}
    %     \begin{shadebox}
    %         右下要素が$0$である行列
    %         \[
    %             A =
    %             \left[
    %                 \begin{array}{ccc|c}
    %                     \multicolumn{3}{c|}{}	& b_1 \\
    %                     \multicolumn{3}{c|}{A'}	& \vdots \\
    %                     \multicolumn{3}{c|}{}	& b_n \\ \hline
    %                     b_1 & \cdots & b_n		& 0
    %                 \end{array}
    %             \right]
    %             =
    %             \left[
    %                 \begin{array}{ccc|c}
    %                     \multicolumn{3}{c|}{} \\
    %                     \multicolumn{3}{c|}{A'}	& \bm{b} \\
    %                     \multicolumn{3}{c|}{} \\ \hline
    %                     \quad & \bm{b}^\top & \quad & 0
    %                 \end{array}
    %             \right]
    %         \]
    %         の行列式は
    %         \[\det{A} = -\bm{b}^\top\tilde{A'}\bm{b} = -\det{(A')}\bm{b}^\top {A'}^{-1}\bm{b}\]
    %         \begin{flushright}
    %             (※$\tilde{A'}$は$A'$の余因子行列)
    %         \end{flushright}
    %     \end{shadebox}
    %     \begin{proof}
    %         \quad\par
    %         第$n$行に関して余因子展開すると
    %         \[\det{A} = \sum_{i=1}^{n}{(-1)^{n+1+i}b_i\det{C_i}}\]
    %         但し$C_i$とは、$[A',\bm{b}]$の第$i$列を除去した小行列である。
    %         $C_i$を第$n$列に関して余因子展開すると
    %         \[\det{C_i} = \sum_{j=1}^{n}{(-1)^{n+j}b_j\det{A'_{ji}}}\]
    %         但し$A'_{ji}$とは、$C_i$の第$n$列と第$j$行を除去したもの、すなわち$A'$の第$j$行と第$i$列を除去した小行列である。
    %         よって
    %         \begin{align*}
    %             \det{A} &= \sum_{i=1}^{n}{(-1)^{n+1+i}b_i\sum_{j=1}^{n}{(-1)^{n+j}b_j\det{A'_{ji}}}} = -\sum_{i=1}^{n}{b_i\sum_{j=1}^{n}{b_j(-1)^{i+j}\det{A'_{ji}}}} \\
    %             &= -\sum_{i=1}^{n}{b_i\sum_{j=1}^{n}{b_j\tilde{a'}_{ij}}} \quad (\tilde{a'}_{ij}\text{は}\tilde{A'}\text{の}(i,j)\text{成分}) \\
    %             &= -\bm{b}^\top\tilde{A'}\bm{b} = -\det{(A')}\bm{b}^\top {A'}^{-1}\bm{b} \quad (\because\;\tilde{A'}=-\det{(A')}{A'}^{-1})
    %         \end{align*}
    %     \end{proof}
    \section{\texorpdfstring{$A\in\realNumbers^{n\times m}, B\in\realNumbers^{m\times n}$}{A∈R\string^(n×m), B∈R\string^(m×n)}に対して\texorpdfstring{$|I_m-AB| = |I_n-BA|$}{|Im-AB| = |In-BA|}}
        \label{I-ABの行列式}
        \begin{proof}
            \[
                M \coloneqq \begin{bmatrix}
                    I_n & A \\
                    B & I_m
                \end{bmatrix}
            \]
            とすると
            \[
                \det{M} = \det{
                    \left(M
                    \begin{bmatrix}
                        I_n & -A \\
                        O & I_m
                    \end{bmatrix}
                    \right)
                } = \det{
                    \begin{bmatrix}
                        I_n & O \\
                        B & I_m-BA
                    \end{bmatrix}
                } = |I_m-BA|
            \]
            一方で
            \[
                \det{M} = \det{
                    \left(M
                    \begin{bmatrix}
                        I_n & O \\
                        -B & I_m
                    \end{bmatrix}
                    \right)
                } = \det{
                    \begin{bmatrix}
                        I_n-AB & A \\
                        ~ & I_m
                    \end{bmatrix}
                } = |I_n-AB|
            \]
        \end{proof}
    \section{系: \texorpdfstring{$\left|I_n-\bm{v}\bm{v}^\top\right| = 1-\norm{\bm{v}}_2^2$}{|In-vv\string^| = 1 - ||v||\string^2}}
        \begin{shadebox}
            $\bm{v} \in \complexNumbers^n$に対して$\left|I_n-\bm{v}\bm{v}^\top\right| = 1-\norm{\bm{v}}_2^2$
        \end{shadebox}
    \section{逆対角転置行列の行列式は元の行列のそれと等しい}
        \begin{proof}
            \quad\par
            $A$を$n$次対角行列とし、その逆対角転置を$A'$とする。
            $S_n$を$n$次の置換全体の集合とする。
            \begin{align*}
                |A'| &= \sum_{\sigma\in S_n} \sgn{\sigma} \prod_{i=1}^n A'[i,\sigma(i)] = \sum_{\sigma\in S_n} \sgn{\sigma} \prod_{i=1}^n A[n-\sigma(i)+1, n-i+1] \\
                &= \sum_{\sigma\in S_n} \sgn{\sigma} \prod_{j=1}^n A[j, n- \sigma^{-1}(n-j+1)+1] \\
                &= \sum_{\sigma\in S_n} \sgn{\sigma} \prod_{j=1}^n A[j, \tilde\sigma(j)] \quad (\tilde\sigma: j \mapsto n+1-\sigma^{-1}(n-j+1))
            \end{align*}
            ここで逆置換の符号が元の置換の符号と等しいこと、および\ref{置換の引数を逆順にしたときの符号}, \ref{出力を逆順にしたときの置換の符号}より$\sgn{\tilde\sigma} = s^2\sgn{\sigma} = \sgn{\sigma}$であるから
            \[ |A'| = \sum_{\sigma\in S_n} \sgn{\tilde\sigma} \prod_{j=1}^n A[j, \tilde\sigma(j)] \]
            $\sigma$と$\tilde\sigma$は一対一で対応するので$\sigma$が$S_n$全体を動くとき$\tilde\sigma$もそうだから上式の$\sum_{\sigma\in S_n}$を$\sum_{\tilde\sigma\in S_n}$で置き換えて
            \[ |A'| = \sum_{\tilde\sigma\in S_n} \sgn{\tilde\sigma} \prod_{j=1}^n A[j, \tilde\sigma(j)] = |A| \]
        \end{proof}